\scanpage{306}
\mspage{5}\deleted{\unsure{5}}
Ce n'est autre que le foncteur représenté par $K_*$, et on le notera aussi
\[
 \deleted{\unsure{K_*(X_*)=}}
 \qquad X_*\longmapsto K_*(X_*)=\operatorname{Hom}(X_*,K_*)
\]
Si $B=(Ab)$, $K_*$ sera de même un objet de $Ss_{ab}$
\deleted{$\unsure{Ab_*}$}, et en effet le foncteur qu'il représente est
automatiquement muni de \deleted{\unsure{structure}}\inserted{\unsure{groupe Abélien}}
\unsure{additive|addition}. $\overline M$ remarque pour les structures
\uline{algébriques} (sur un ou plusieurs objets de base) qui
``puissent s'exprimer en termes de $\varprojlim$ exclusivement''
(gps, \deleted{\unsure{anneaux}}, anneaux, modules sur
\deleted{\unsure{des anneaux}}\deleted{\unsure{etc.}} etc.).

\uline{Rem}. Les \deleted{foncteurs} invariants ainsi obtenus,
à coups de foncteurs représent.bles --- plus génér., d'objets simpliciaux
ou cosimpliciaux de catégories $B$ ou $A$ --- sont de nature trop
``grosse'' et ``fruste'' pour être intéressants directement --- en
particulier, ce ne sont pas des invariants du type d'homotopie de $X_*$.
On devra en extraire des invariants plus subtils qui soient des invariants
du type d'homotopie de $X_*$, --- ceux-ci ne seront pas de nature si simple
--- i.e. ``repr.'' par des objets simpliciaux ou cosimpliciaux.
Notre propos sera de mettre sur les invariants ``gros'' qu'on vient de décrire,
un maximum de structures supplémentaires remarquables, et dans les catégories
de structures de ce type (complexes de chaînes ou de cochaînes, algèbres
cosimpliciales etc.) d'inverser les flèches qu'on obtient à partir d'équiv.
d'homotopie des $Ss$, et de passer à des ``catégories de fractions'' en
inversant ces flèches. L'invariant ainsi obtenu
\deleted{\unsure{deviendra}}\inserted{(plus ``grossier''}
bien sûr, mais plus subtil dans sa définition, et moins ``redondant'')
\unsure{exprimera|exprime} donc de façon plus ou moins complète le type
d'homotopie, ou en saisira de façon adéquate tels ou tels aspects.

Dans le cas où $A$ est une catégorie additive karoubienne (tout projecteur
dans un objet a un noyau, ou ce qui revient au même, tt proj. a un conoyau)
le th.\unsure{ème|éorème} de Dold-Puppe donne des équivalences canoniques
\diagram{\begin{tikzcd}[column sep=large]
 A_*\arrow[r,shift left=1ex,"{ND_{\cdot}}","\approx"']&
 A_{\cdot}\arrow[l,shift left=1ex,"{DP_*}","\approx"']
\end{tikzcd}},
\margin{$ND=$ ``non-dégénéré''\\$DP=$ ``Dold-Puppe''}
où $A_{\cdot}$ désigne la catégorie des \uline{complexes de chaînes}
dans $A$. On désigne par une même lettre \inserted{telle $K$} deux objets
qui se correspondent, mais en mettant en indice soit $*$ soit $\cdot$,
suivant qu'il
\margin{307}
